\documentclass{article}
\usepackage{amsmath,amsthm,amssymb}
\usepackage{mathtext}
\usepackage[T2A]{fontenc}
\usepackage[utf8]{inputenc}
\usepackage[english,russian]{babel}
\usepackage{graphicx}
\usepackage{booktabs}
\usepackage{tikz}
\usepackage{hyperref}

\title{Argument Mining для кинорецензий:\\воспроизведение бейзлайна и диагностика утечки оценки}
\author{Анастасия Шпилева}
\date{Сентябрь 2026}

\begin{document}
\maketitle

\begin{abstract}
Мы воспроизводим бейзлайн из работы <<Argument Mining for Film Reviews>>~\cite{serbina2026argument}:
энкодер для русского языка с двумя классификационными головами, предсказывающий по тексту
кинорецензии позицию автора (stance) и полярность его аргументации (premise). Помимо
воспроизведения мы показываем, что в 56\% текстов датасета лежит явная числовая оценка фильма
(<<8 из 10>>) --- ровно тот балл, из которого выведен таргет stance, --- и измеряем, сколько
качества на ней держится. Удаление оценки стоит rubert-base-cased $7.78 \pm 2.00$ пункта
macro-F1 на stance и почти ничего не стоит ruRoberta-large, из-за чего исходная постановка
занижает разницу между энкодерами примерно вчетверо. Отдельно мы показываем, что разброс
по seed'ам обучения ($\pm 2.43$ пункта) сопоставим с измеряемыми эффектами, так что
одиночный прогон в этой задаче не интерпретируем в принципе. Код:
\url{https://github.com/banka-lecho/nlp_project}.
\end{abstract}

\section{Введение}

Argument mining --- извлечение из текста не только того, \emph{что} человек думает, но и
\emph{почему}: какие посылки он приводит и в какую сторону они работают. Для отзывов
это содержательное отличие от обычного сентимент-анализа: рецензент может поставить фильму
высокий балл, разругав его по всем пунктам, или наоборот --- выставить низкую оценку,
перечислив достоинства. В датасете, с которым мы работаем, такие неконгруэнтные отзывы
составляют 29.3\%, и именно на них разваливается качество всех моделей (раздел~\ref{sec:results}).

Работа~\cite{serbina2026argument} предлагает для этой задачи простой и воспроизводимый бейзлайн:
один энкодер, две линейные головы, две эпохи обучения. Наша работа устроена в два слоя.

\textbf{Первый слой --- воспроизведение.} В статье не зафиксированы learning rate, размер
батча, максимальная длина и seed разбиения, поэтому точное повторение чисел невозможно;
мы берём стандартные значения, повторяем протокол и сравниваем результаты через
доверительные интервалы.

\textbf{Второй слой --- диагностика.} Разведочный анализ данных показал, что таргет stance
выведен из колонки \texttt{grade10}, а сама эта оценка в 56\% случаев дословно присутствует
в тексте рецензии (<<7/10>>, <<8 баллов из 10>>) и в 98.5\% таких случаев --- в последних
15\% текста. То есть модель может не читать аргументацию вовсе, а находить в конце число
и переводить его в класс. Мы проверяем эту гипотезу тремя независимыми способами:
правилом без всякой модели, дообучением на редактированных текстах и линейным пробингом
замороженных представлений.

Ключевое отличие от исходной постановки --- в том, что каждое число здесь сопровождается
оценкой его собственной неустойчивости: по seed'ам обучения, по перезапуску с тем же seed'ом
и по бутстрэпу тестовой выборки. Как выяснилось, для половины интересных сравнений эта
неустойчивость и есть главный результат.

\subsection{Состав работы}

\textbf{Анастасия Шпилева} выполнила всю работу: разведочный анализ данных, реализацию
бейзлайна и всех расширений, прогоны, анализ результатов и текст отчёта.

\section{Обзор существующих подходов}
\label{sec:related}

\paragraph{Argument mining.} Классическая постановка задачи --- выделение в тексте
аргументативных единиц и связей между ними~\cite{mochales2011argumentation,lippi2016argumentation};
эталонный корпус --- persuasive essays со схемой <<major claim / claim / premise>> и
разметкой поддержки и атаки~\cite{stab2017parsing}. Для русского языка ближайший ориентир ---
дорожка RuArg-2022~\cite{kotelnikov2022ruarg}, где задача тоже разложена на две параллельные
классификации: позиция автора по отношению к утверждению и наличие/полярность посылки.
Постановка~\cite{serbina2026argument} наследует эту двухголовую структуру, но переносит её
с твитов о ковиде на кинорецензии, где позиция дополнительно закреплена числовой оценкой.

\paragraph{Энкодеры для русского.} Стандартный базовый энкодер --- rubert-base-cased,
получен дообучением многоязычного BERT~\cite{devlin2019bert} на русских данных с
пересобранным словарём~\cite{kuratov2019adaptation}. Более сильная модель ---
ruRoberta-large из семейства предобученных моделей SberDevices~\cite{zmitrovich2024family},
воспроизводящая рецепт RoBERTa~\cite{liu2019roberta}. Обе модели используются
в~\cite{serbina2026argument} и обе взяты нами.

\paragraph{Артефакты разметки и поверхностные признаки.} Линия работ, к которой примыкает
наша диагностическая часть: модель показывает высокое качество, пользуясь признаком,
не имеющим отношения к декларируемой задаче. В SNLI гипотеза сама по себе, без посылки,
предсказывает метку сильно выше случайного~\cite{gururangan2018annotation,poliak2018hypothesis};
в ряде QA-бенчмарков вопрос можно не читать~\cite{kaushik2018much}; в Argument Reasoning
Comprehension результат BERT целиком объяснялся частотностью отдельных
слов-подсказок~\cite{niven2019probing}. Методологически стандартный ответ --- ablation
входа: удалить подозрительную часть текста и посмотреть, сколько качества осталось.
Мы применяем ровно эту схему, но делаем её \emph{накопительной} (лестница из пяти ступеней),
чтобы отделить потерю утечки от потери собственно аргументации.

\paragraph{Устойчивость результатов.} Дообучение BERT-подобных моделей на выборках
в несколько тысяч примеров сильно зависит от инициализации и порядка данных: разброс
по seed'ам может превышать разницу между сравниваемыми
методами~\cite{dodge2020fine,reimers2017reporting}. Отсюда же требование не публиковать
одиночное число~\cite{reimers2017reporting} и осторожность со стандартными
разбиениями~\cite{gorman2019we}. Для сравнения двух систем на одной тестовой выборке
применяется парный бутстрэп~\cite{koehn2004statistical,efron1993introduction};
общий обзор процедур значимости в NLP --- \cite{dror2018hitchhiker}.

\paragraph{Порядковая структура классов и калибровка.} Классы Bad $<$ Neutral $<$ Good
упорядочены, и ошибка Bad$\to$Good хуже ошибки Bad$\to$Neutral; номинальный softmax этого
не знает. Классические ответы --- задачи rating inference~\cite{pang2005seeing},
ранг-согласованные головы CORAL~\cite{cao2020rank} и CORN~\cite{shi2021deep}, метрика
quadratic weighted kappa~\cite{cohen1968weighted}. Уверенность современных сетей
систематически завышена, и стандартная поправка --- температурное
шкалирование по валидации~\cite{guo2017calibration}.

\paragraph{Пробинг.} Линейный классификатор поверх замороженных представлений --- дешёвый
способ узнать, какая информация лежит в них <<на поверхности>>, до всякого
дообучения~\cite{alain2016understanding,tenney2019bert}. Мы используем пробинг не для
анализа слоёв, а как дешёвую замену дорогому ablation-эксперименту: вся лестница из пяти
ступеней для обеих моделей считается за 41 минуту против примерно десяти часов дообучения.

\paragraph{Результаты исходной работы.} Числа, с которыми мы сравниваемся, ---
macro-F1 на общей тестовой выборке (Table 4, General в~\cite{serbina2026argument}):
rubert-base-cased 71.31 / 75.00 и ruRoberta-large 75.56 / 78.10 для stance и premise
соответственно.

\section{Описание подхода}

\subsection{Бейзлайн}

Формально: по тексту рецензии $x$ предсказываются две метки из одного и того же
упорядоченного множества $\{\text{Bad}, \text{Neutral}, \text{Good}\}$ ---
$y^{\text{st}}$ (позиция автора, колонка \texttt{NEW\_grade3}) и $y^{\text{pr}}$
(полярность аргументации, колонка \texttt{arg\_label}).

Энкодер общий, головы независимы (Рис.~\ref{fig:model}). Пусть $h(x) \in \mathbb{R}^d$ ---
представление \texttt{[CLS]} после dropout, тогда
$$z^{t} = W^{t} h(x) + b^{t}, \qquad t \in \{\text{st}, \text{pr}\},$$
а лосс --- сумма двух кросс-энтропий, без взвешивания задач:
$$\mathcal{L} = \mathrm{CE}(z^{\text{st}}, y^{\text{st}}) + \mathrm{CE}(z^{\text{pr}}, y^{\text{pr}}).$$
Обучаются все параметры, включая энкодер, две эпохи.

\begin{figure}[!tbh]
    \centering
    \begin{tikzpicture}[x=1mm,y=1mm,font=\small]
      \draw (0,0) rectangle (46,8); \node at (23,4) {текст рецензии $x$};
      \draw[->] (23,8) -- (23,14);
      \draw (0,14) rectangle (46,24); \node at (23,19) {энкодер (rubert / ruRoberta)};
      \draw[->] (23,24) -- (23,30);
      \draw (10,30) rectangle (36,38); \node at (23,34) {\texttt{[CLS]} + dropout};
      \draw[->] (17,38) -- (9,44); \draw[->] (29,38) -- (37,44);
      \draw (-4,44) rectangle (22,52); \node at (9,48) {голова stance};
      \draw (24,44) rectangle (50,52); \node at (37,48) {голова premise};
      \draw[->] (9,52) -- (9,58); \draw[->] (37,52) -- (37,58);
      \node at (9,61) {$\mathrm{CE}$}; \node at (37,61) {$\mathrm{CE}$};
      \draw[->] (12,61) -- (20,61); \draw[->] (34,61) -- (26,61);
      \node at (23,61) {$+$};
    \end{tikzpicture}
    \caption{Архитектура бейзлайна: общий энкодер, две параллельные линейные головы,
    лосс --- сумма двух кросс-энтропий.}
    \label{fig:model}
\end{figure}

\subsection{Лестница удаления текста}
\label{sec:ladder}

Чтобы измерить вклад явной оценки, мы обучаем и оцениваем те же модели на текстах
с последовательно вырезаемыми фрагментами. Лестница \emph{накопительная}: каждая
следующая ступень убирает строго больше предыдущей, и начиная со второй ступени явной
оценки не остаётся ни в одном тексте (проверено: 0 остаточных вхождений).

\begin{table}[!tbh]
\centering
\begin{tabular}{llr}
\toprule
Ступень & Что остаётся в тексте & Слов в среднем \\
\midrule
\texttt{none}          & всё                                        & 103.5 \\
\texttt{grade}         & текст без числовой оценки                  & 102.2 \\
\texttt{last-sentence} & без оценки и без последнего предложения     & 86.6 \\
\texttt{tail15}        & без оценки и без последних 15\% символов    & 86.4 \\
\texttt{first-half}    & без оценки, только первая половина          & 50.5 \\
\bottomrule
\end{tabular}
\caption{Лестница удаления (\texttt{analysis/redaction\_ladder.csv}). Оценка снимается
регулярными выражениями, покрывающими варианты <<8 из 10>>, <<8/10>>, <<8 баллов из десяти>>
и т.\,п.; в среднем удаляется 1.28 слова.}
\label{tab:ladder}
\end{table}

Смысл накопительности: падение на ступени \texttt{grade} --- это потеря утечки, а падение
\emph{ниже} этого уровня уже нельзя списать на балл, это потеря собственно аргументации.

\subsection{Линейный пробинг}

Энкодер замораживается, тексты кодируются усреднением по токенам, признаки стандартизуются,
поверх учится логистическая регрессия; коэффициент регуляризации $C$ подбирается на валидации
(иначе пробинг меряет не представления, а удачный выбор регуляризации). Прогоняется вся
лестница для обеих моделей.

\subsection{Расширения сверх статьи}

Все перечисленное реализовано и по умолчанию выключено, так что команда по умолчанию
воспроизводит ровно постановку бейзлайна.

\textbf{Ординальная голова.} \texttt{--head corn} заменяет три независимых логита на две
условные вероятности $P(y > \text{Bad})$ и $P(y > \text{Neutral} \mid y > \text{Bad})$
по схеме CORN~\cite{shi2021deep}: узел $k$ учится только на подвыборке с $y > k-1$,
что даёт ранг-согласованные вероятности
$P(y = k) = \prod_{j<k}\sigma(z_j)\,(1 - \sigma(z_k))$.

\textbf{Веса классов.} Метрика --- macro-F1, а лосс --- невзвешенный CE, при том что самый
редкий класс и идёт хуже всех. \texttt{--class-weights balanced} согласует лосс с метрикой.

\textbf{Групповой сплит.} \texttt{--group movie\_name} гарантирует, что ни один фильм
не попадает одновременно в train и test (в исходном сплите 95.5\% фильмов теста есть
в трейне, так что обобщение на новый фильм не измеряется~\cite{gorman2019we}).

\textbf{Отбор эпохи, пулинг, калибровка.} \texttt{--select-by mean} берёт эпоху с лучшим
средним val macro-F1 вместо последней; \texttt{--pooling mean} --- усреднение вместо
\texttt{[CLS]}; \texttt{evaluate.py --calibrate} подбирает температуру на валидации
и применяет к тесту~\cite{guo2017calibration}.

\subsection{Конкурирующие подходы и бейзлайны}
\label{sec:baselines}

\textbf{Правило по числу в тексте (без обучения).} Из текста извлекается число перед
<<из 10>>, и метка ставится по порогам: $\le 4 \to$ Bad, $\le 6 \to$ Neutral, иначе Good.
Правило применимо к 1818 отзывам, где оценка присутствует. Это прямой аналог
hypothesis-only бейзлайна~\cite{poliak2018hypothesis}: оно не читает текст вообще.

\textbf{Линейный пробинг.} Нижняя граница для дообучения: сколько сигнала доступно
без всякой адаптации энкодера.

\textbf{Числа из статьи.} Верхний ориентир для воспроизведения.

\section{Данные}

\subsection{Описание}

Датасет --- \texttt{otipl2125/film\_review\_argumentation}\footnote{\url{https://huggingface.co/datasets/otipl2125/film_review_argumentation},
лицензия CC-BY-4.0, доступен для исследовательских целей.}, представленный
в~\cite{serbina2026argument}. Это 3280 русскоязычных рецензий на 323 фильма от 2666 авторов,
собранных с кинопортала: 2226 отзывов на фильмы из <<топ-250>> и 1054 --- на фильмы из
<<худшей сотни>>; даты публикации --- с июля 2004 по ноябрь 2012 года.

Разметка двухслойная. Колонка \texttt{grade3} --- трёхклассовая свёртка авторской оценки
\texttt{grade10}; колонка \texttt{NEW\_grade3} (таргет stance) --- та же свёртка после
ручной проверки, она отличается от автоматической в 12.5\% случаев. Колонка
\texttt{arg\_label} (таргет premise) --- размеченная полярность аргументации.
Совместное распределение меток приведено в Таб.~\ref{tab:joint}: конгруэнтны, то есть
совпадают, 70.7\% пар, а крайние расхождения (Bad$\leftrightarrow$Good) редки ---
79 пар против 457 соседних, что и мотивирует ординальную голову.

\begin{table}[!tbh]
\begin{center}
\begin{tabular}[t]{|l|ccc|c|}
\hline
stance \textbackslash{} premise & Bad & Neutral & Good & всего \\
\hline
Bad     & 675 & 398 & 47   & 1120 \\
Neutral & 59  & 514 & 178  & 751  \\
Good    & 32  & 247 & 1130 & 1409 \\
\hline
всего   & 766 & 1159 & 1355 & 3280 \\
\hline
\end{tabular}
\caption{Совместное распределение меток. По строкам --- stance (\texttt{NEW\_grade3}),
по столбцам --- premise (\texttt{arg\_label}). Доли классов: stance 34/23/43\%,
premise 23/35/41\%.}
\label{tab:joint}
\end{center}
\end{table}

Разбиение --- 70/15/15 со стратификацией по \texttt{arg\_label}, как в статье:
2296 / 492 / 492 (Таб.~\ref{tab:statistics}). Распределение stance при этом отдельно
не контролируется и держится на seed 42 случайно.

\begin{table}[tbh!]
\begin{center}
\begin{tabular}[t]{|l|ccc|}
\hline
 & Train & Valid & Test \\
\hline
Отзывов & 2296 & 492 & 492 \\
stance Bad/Neutral/Good & .344/.226/.430 & .333/.234/.433 & .339/.236/.425 \\
premise Bad/Neutral/Good & .233/.353/.413 & .234/.354/.413 & .234/.354/.413 \\
Слов в отзыве & \multicolumn{3}{c|}{$103.0 \pm 36.0$ (min 5, max 156)} \\
Токенов (max) & \multicolumn{3}{c|}{305 rubert / 356 ruRoberta} \\
Дубликатов текста & \multicolumn{3}{c|}{0} \\
\hline
\end{tabular}
\caption{Статистика датасета и разбиения (\texttt{analysis/splits.csv},
\texttt{analysis/lengths.csv}). При \texttt{max\_length}=512 обрезка не срабатывает
ни разу ни для одной модели.}
\label{tab:statistics}
\end{center}
\end{table}

\subsection{Что показал разведочный анализ}
\label{sec:eda}

Ниже --- находки, влияющие на интерпретацию результатов; полный разбор в
\texttt{src/analyze.py}, таблицы --- в \texttt{analysis/}.

\textbf{Явная оценка в тексте.} 56\% отзывов содержат конструкцию вида <<8 из 10>>,
<<10 баллов из 10>>, <<7/10>>, и у 98.5\% из них она стоит в последних 15\% текста.
Извлечённое число в 92.8\% случаев равно колонке \texttt{grade10}, \emph{из которой
выведен таргет} \texttt{NEW\_grade3}. Иными словами, в поле \texttt{content} лежит
сам источник метки.

\textbf{Насколько этого достаточно.} Правило из раздела~\ref{sec:baselines} даёт
на 1818 отзывах с оценкой macro-F1 0.718 на stance и 0.647 на premise
(accuracy 0.803 и 0.687) --- без единого обученного параметра и при том, что
дообученная rubert-base-cased даёт 72.20 на полном тесте.

\textbf{Шум.} 32\% текстов приходят с пробельным паддингом по краям (снимается
\texttt{strip}), 59.7\% содержат склейки предложений вида \texttt{Темнота.Чужие},
1.3\% --- двойные пробелы внутри. HTML-тегов, URL и буквы <<ё>> нет.

\textbf{Сплит не групповой.} 95.5\% фильмов из теста встречаются в трейне; у одного
автора до 18 отзывов. Группировка по автору почти ничего не даёт (медиана --- один
отзыв на автора), группировка по фильму убирает пересечение полностью.

\textbf{Не баг.} 452 неуникальных \texttt{review\_id} --- это переиспользованный
идентификатор у разных отзывов о разных фильмах (307 id встречаются больше чем
у одного фильма); тексты при этом уникальны, дедуплицировать по \texttt{review\_id} нельзя.

\section{Эксперименты}

\subsection{Метрики}

Основная метрика --- macro-F1, как в~\cite{serbina2026argument}:
$$F_1^{\text{macro}} = \frac{1}{K}\sum_{k=1}^{K} \frac{2\,P_k R_k}{P_k + R_k},
\qquad K = 3.$$
Доверительные интервалы --- бутстрэп по тестовой выборке, 1000 итераций,
перцентильный 95\% ДИ~\cite{efron1993introduction}, отдельно по подгруппам
(век выхода фильма, рейтинговая часть, конгруэнтность меток).

Для сравнения двух прогонов на \emph{одной} тестовой выборке пересечение отдельных ДИ
неинформативно: ошибки моделей скоррелированы. Используется парный
бутстрэп~\cite{koehn2004statistical}: на каждой итерации ресэмплится один и тот же набор
индексов для обеих систем, и оценивается распределение разности
$\Delta = F_1(A) - F_1(B)$; $p$-value --- доля итераций с $\Delta \le 0$.

Дополнительно считаются метрики, учитывающие порядок классов:
quadratic weighted kappa~\cite{cohen1968weighted}
$$\kappa = 1 - \frac{\sum_{i,j} w_{ij} O_{ij}}{\sum_{i,j} w_{ij} E_{ij}},
\qquad w_{ij} = \frac{(i-j)^2}{(K-1)^2},$$
MAE по индексам классов и доля <<ошибок через класс>> ($|i - j| = 2$).
Для калибровки --- ECE (взвешенный разрыв между уверенностью и точностью по 15 бинам)
и NLL~\cite{guo2017calibration}.

\subsection{Постановка эксперимента}

Гиперпараметры, не указанные в статье, взяты стандартные: lr $2 \cdot 10^{-5}$,
batch 16, \texttt{max\_length} 512, warmup 10\%, weight decay 0.01, AdamW, клиппинг
градиента 1.0, 2 эпохи, берётся последняя эпоха. ruRoberta-large обучалась с
\texttt{--batch-size 4 --grad-accum 4}: на 24\,GB unified memory батч 16 уходит в своп.
Железо --- Mac M4 (MPS); rubert-base --- 11 минут на прогон, ruRoberta-large --- 34 минуты.

Сплит фиксирован (\texttt{--split-seed 42}) во всех прогонах, меняется только seed
инициализации --- так тестовая выборка у всех прогонов общая и парный бутстрэп остаётся
применим. rubert-base-cased прогонялась на 5 seed'ах (42--46) в двух режимах
(\texttt{none} и \texttt{grade}), ruRoberta-large --- пока на одном seed в тех же двух
режимах. Разница по seed'ам считается попарно (одинаковая инициализация против
одинаковой), а не как разность двух средних.

Поскольку seed разбиения в статье не указан, состав нашей тестовой выборки заведомо
отличается от авторского, и корректно сравнивать результаты по доверительным интервалам,
а не по совпадению чисел.

\section{Результаты}
\label{sec:results}

\subsection{Воспроизведение}

\begin{table}[!tbh]
\centering
\begin{tabular}{lcccc}
\toprule
Модель & Stance & Premise & Stance~\cite{serbina2026argument} & Premise~\cite{serbina2026argument} \\
\midrule
rubert-base-cased, 5 seeds & $72.20 \pm 2.43$ & $69.09 \pm 1.32$ & 71.31 & 75.00 \\
ruRoberta-large, 1 seed    & 74.22            & 73.66            & 75.56 & 78.10 \\
\bottomrule
\end{tabular}
\caption{macro-F1 на тесте ($n=492$), среднее $\pm$ стандартное отклонение по seed'ам
обучения. Сплит у всех прогонов один.}
\label{tab:reproduction}
\end{table}

Stance воспроизводится: 72.20 против 71.31 у авторов, разница внутри собственного
разброса. Premise стабильно ниже на 4--6 пунктов у обеих моделей --- вероятное объяснение
в том, что неуказанные гиперпараметры (прежде всего lr и число эпох) подобраны авторами
под эту задачу, а мы взяли одни значения на обе.

Главное наблюдение здесь --- не среднее, а разброс. У rubert-base stance по пяти seed'ам
укладывается в диапазон 68.86--75.00, размах 6.14 пункта. Опубликовать из этого диапазона
одно число и сравнить его с 71.31 из статьи бессмысленно: с равной вероятностью получилось
бы <<на 3.7 хуже>> и <<на 3.7 лучше>>. Premise устойчивее: 67.05--70.03, размах 2.98.
Это ровно то, о чём предупреждают~\cite{dodge2020fine,reimers2017reporting}.

\subsection{Прогон невоспроизводим побитово даже при фиксированном seed}

Инициализация детерминирована (веса голов совпадают побитово, первые три шага дают
одинаковый лосс до восьмого знака), но с четвёртого шага лоссы расходятся на $\sim 10^{-7}$:
на MPS порядок суммирования в редукциях не гарантирован. За 144 шага $\times$ 2 эпохи это
накапливается примерно в пункт macro-F1. На CUDA то же самое по умолчанию, но там лечится
\texttt{torch.use\_deterministic\_algorithms(True)}; на MPS такого рычага нет.
Один и тот же эксперимент, тот же seed 42, дважды:

\begin{table}[!tbh]
\centering
\begin{tabular}{lll}
\toprule
Задача & Прогон 1 & Прогон 2 \\
\midrule
stance  & $+6.85$, ДИ $(+2.61, +10.95)$, $p=0.001$ & $+8.52$, ДИ $(+4.39, +12.68)$, $p=0.000$ \\
premise & $+0.47$, ДИ $(-3.32, +4.10)$, $p=0.404$  & $+2.71$, ДИ $(-0.78, +6.25)$, $p=0.066$ \\
\bottomrule
\end{tabular}
\caption{Дельта <<с оценкой минус без оценки>> для rubert-base-cased, seed 42, два запуска
одной и той же команды.}
\label{tab:nondet}
\end{table}

Вывод по stance устойчив, а по premise <<уверенно ничего>> ($p=0.40$) превратилось
в <<на границе>> ($p=0.066$) от простого перезапуска. Бутстрэп-ДИ меряет шум выборки,
но у точечной оценки есть собственный шум перезапуска сопоставимой величины,
и одиночный ДИ его не видит.

\subsection{Сколько качества держится на явной оценке}

\begin{table}[!tbh]
\centering
\begin{tabular}{lcc}
\toprule
Модель & Stance & Premise \\
\midrule
rubert-base-cased, 5 seeds & $64.42 \pm 1.15$ & $68.09 \pm 0.79$ \\
ruRoberta-large, 1 seed    & 72.93            & 74.27 \\
\bottomrule
\end{tabular}
\caption{macro-F1 на тесте после удаления явной оценки (\texttt{--redact grade}).}
\label{tab:nograde}
\end{table}

\begin{table}[!tbh]
\centering
\begin{tabular}{llccl}
\toprule
Модель & Задача & $\Delta$ macro-F1 & Разброс & По seed'ам \\
\midrule
rubert-base, 5 seeds & stance  & $\mathbf{+7.78}$ & $\pm 2.00$ & $+8.52\ +6.36\ +10.52\ +5.36\ +8.15$ \\
rubert-base, 5 seeds & premise & $+1.01$ & $\pm 1.16$ & $+2.71\ +0.81\ -0.53\ +1.20\ +0.84$ \\
ruRoberta, 1 seed    & stance  & $+1.29$ & ДИ $(-1.78, +4.52)$ & $p=0.208$ \\
ruRoberta, 1 seed    & premise & $-0.62$ & ДИ $(-3.45, +2.17)$ & $p=0.655$ \\
\bottomrule
\end{tabular}
\caption{Потеря качества при удалении оценки (бейзлайн минус прогон без оценки).
Для rubert --- попарно по совпадающим seed'ам, для ruRoberta --- парный бутстрэп
единственной пары прогонов.}
\label{tab:delta}
\end{table}

У rubert-base на stance дельта положительна на всех пяти seed'ах и вчетверо больше
собственного разброса --- вывод устойчив: база держалась именно на числе в конце отзыва.
На premise дельта $+1.01$ при разбросе $\pm 1.16$ и одном отрицательном значении из пяти,
то есть от нуля неотличима. ruRoberta-large не теряет ничего ни на одной задаче.

Это ожидаемо в части задач: для \texttt{arg\_label} балл предсказывает хуже
(macro-F1 правила 0.647 против 0.718), поэтому на premise не теряет никто. Но расхождение
\emph{между моделями} на stance --- содержательный результат: сильный энкодер читает
аргументацию, слабый --- финальную цифру.

Побочное наблюдение: без оценки в тексте разброс по seed'ам падает вдвое
($2.43 \to 1.15$ на stance). Похоже, часть нестабильности бейзлайна и есть утечка:
в зависимости от инициализации модель цепляется за балл в разной степени.

\subsection{Утечка искажает сравнение моделей}

\begin{table}[!tbh]
\centering
\begin{tabular}{llccc}
\toprule
Данные & Задача & ruRoberta $-$ rubert & Разброс rubert & Разрыв в $\sigma$ \\
\midrule
бейзлайн   & stance  & $+2.01$ & $\pm 2.43$ & 0.8 \\
без оценки & stance  & $\mathbf{+8.51}$ & $\pm 1.15$ & $\mathbf{7.4}$ \\
бейзлайн   & premise & $+4.56$ & $\pm 1.32$ & 3.5 \\
без оценки & premise & $+6.18$ & $\pm 0.79$ & 7.8 \\
\bottomrule
\end{tabular}
\caption{Разрыв между энкодерами до и после удаления утечки. ruRoberta (1 seed)
сравнивается со средним rubert по 5 seed'ам, разрыв меряется в единицах разброса rubert.}
\label{tab:gap}
\end{table}

На исходных данных преимущество ruRoberta по stance тонет в разбросе rubert: 2 пункта
против $\sigma = 2.43$. Цифра в конце отзыва доступна обеим моделям и подтягивает слабую
до уровня сильной. После удаления утечки разрыв вырастает вчетверо, а разброс слабой модели
вдвое падает, и отрыв становится семикратным по $\sigma$. То есть бейзлайн из статьи
\emph{недооценивает} разницу между энкодерами ровно там, где она есть.

Оговорка: у ruRoberta один seed, её собственная неопределённость в эти числа не входит,
так что две последние колонки --- верхняя оценка уверенности.

\subsection{Линейный пробинг по лестнице}

\begin{table}[!tbh]
\centering
\begin{tabular}{lcccc}
\toprule
Ступень & rubert stance & rubert premise & ruRoberta stance & ruRoberta premise \\
\midrule
\texttt{none}          & 60.66 & 65.62 & 67.48 & 69.69 \\
\texttt{grade}         & 56.31 ($\mathbf{-4.35}$) & 64.11 ($-1.51$) & 65.77 ($-1.71$) & 70.22 ($+0.53$) \\
\texttt{last-sentence} & 55.83 ($-4.83$) & 62.16 ($-3.46$) & 62.57 ($-4.91$) & 68.52 ($-1.17$) \\
\texttt{tail15}        & 56.05 ($-4.61$) & 61.99 ($-3.63$) & 62.95 ($-4.53$) & 69.68 ($-0.01$) \\
\texttt{first-half}    & 51.52 ($-9.14$) & 58.52 ($-7.10$) & 58.14 ($-9.34$) & 65.39 ($-4.30$) \\
\bottomrule
\end{tabular}
\caption{macro-F1 на тесте, замороженный энкодер плюс логистическая регрессия;
в скобках --- изменение относительно \texttt{none} (\texttt{analysis/probe\_*.csv}).}
\label{tab:probe}
\end{table}

Здесь видно ровно то же расхождение между моделями, что и при дообучении, но уже
\emph{без} дообучения: на первой ступени rubert теряет на stance 4.35 пункта,
а ruRoberta --- 1.71. У дообученных моделей та же пара --- $-6.85$ ($p=0.001$)
и $+1.29$ (незначимо). Асимметрия не создаётся дообучением: она уже лежит
в представлениях, дообучение её только усиливает.

Дальше по лестнице модели сходятся: к \texttt{first-half} обе теряют около 9 пунктов
на stance. В этом и смысл накопительности --- ниже уровня \texttt{grade} оценки в тексте
нет ни у кого, и одинаковое падение у обеих моделей означает, что теряется общая
для них аргументация, а не утечка.

Отдельно стоит поведение задач. stance у rubert проседает на одном удалённом числе
и дальше выходит на плато: убрать ещё 16 слов хвоста стоит примерно ноль. premise так
себя не ведёт ни у одной модели --- он деградирует монотонно вместе с длиной текста.
Сигнал stance сосредоточен в балле, сигнал premise размазан по аргументации.

\subsection{Подгруппы}

\begin{table}[!tbh]
\centering
\begin{tabular}{lrcccc}
\toprule
& & \multicolumn{2}{c}{rubert-base-cased} & \multicolumn{2}{c}{ruRoberta-large} \\
\cmidrule(lr){3-4}\cmidrule(lr){5-6}
Подгруппа & $n$ & stance & premise & stance & premise \\
\midrule
General       & 492 & 73.72 & 70.03 & 74.22 & 73.66 \\
XX век        & 145 & 65.92 & 63.51 & 72.09 & 79.20 \\
XXI век       & 347 & 74.14 & 70.21 & 73.32 & 72.05 \\
топ-250       & 325 & 71.69 & 66.95 & 72.39 & 73.68 \\
худшая сотня  & 167 & 61.57 & 62.13 & 59.61 & 66.24 \\
Конгруэнтные  & 357 & 77.29 & 76.73 & 80.82 & 81.39 \\
Неконгруэнтные & 135 & \textbf{61.18} & \textbf{45.98} & \textbf{52.95} & \textbf{50.04} \\
\bottomrule
\end{tabular}
\caption{macro-F1 по подгруппам, seed 42 (\texttt{runs/*/report.csv}; полные отчёты
содержат также 95\% ДИ, QWK, MAE и долю ошибок через класс).}
\label{tab:subgroups}
\end{table}

Разрыв между конгруэнтными и неконгруэнтными отзывами --- 16 пунктов у rubert и 28--31
у ruRoberta. На отзывах, где позиция и аргументация расходятся, premise падает до
уровня 46--50 macro-F1, то есть модель фактически предсказывает аргументацию по позиции.
Это и есть содержательно интересная часть задачи, и она решается хуже всего.
Вторая просевшая подгруппа --- фильмы из худшей сотни (60--62 на stance).

\subsection{Калибровка}

Температура, подобранная на валидации: $T = 1.15$ для stance и $1.23$ для premise
(rubert, seed 42), то есть модель переуверена умеренно. ECE снижается с 0.043 до 0.034
(stance) и с 0.075 до 0.070 (premise), NLL --- незначительно. macro-F1 при этом не меняется:
argmax инвариантен к температуре.

\subsection{Пример работы модели}

\begin{table}[!tbh]
    \centering
    \begin{tabular}{|p{12cm}|}
\hline
\textbf{Отзыв на фильм <<Ведьма (2006)>>, тестовая выборка:}\newline
<<Это первый российский ужастик. Упор в фильме сделан на атмосферу и нагнетение
обстановки, а не на тупую резню, коей славяться тупые американские ужастики
(не все но многие).Фильм получился качественный в котором сочитается сюжет,
глубокая идея, пугающая и напряженная атмосфера, хорошие специфекты и даже
немного юмора.\underline{7 из 10}>> \\
\hline
\textbf{Истина:} stance = Neutral, premise = Good \\
\textbf{Бейзлайн (\texttt{none}):} stance = Neutral, premise = \textbf{Neutral} \\
\textbf{Без оценки (\texttt{grade}):} stance = Neutral, premise = \textbf{Good} \\
\hline
    \end{tabular}
    \caption{Реальный пример из теста (rubert-base-cased, seed 42). Балл 7/10 даёт
    stance = Neutral, и бейзлайн переносит этот же класс на вторую голову, хотя текст
    состоит из одних похвал. Без балла модель вынуждена читать аргументацию и ставит
    premise = Good правильно. Орфография источника сохранена; подчёркивание --- наше,
    это фрагмент, который вырезает \texttt{--redact grade}.}
    \label{tab:output}
\end{table}

Такая картина типична: среди 492 тестовых отзывов 13 --- случаи, где прогон без оценки
исправляет premise, и почти все они неконгруэнтные, с баллом в последнем предложении.

\section{Что сделано и что осталось}

Сделано: воспроизведён бейзлайн из~\cite{serbina2026argument} (общий энкодер, две головы,
2 эпохи) на двух русских энкодерах; проведён разведочный анализ, обнаруживший утечку
таргета stance в текст; реализованы и посчитаны три независимых способа её измерить
(правило без обучения, накопительная лестница удаления с дообучением, линейный пробинг);
собран статистический аппарат --- бутстрэп по подгруппам, парный бутстрэп для сравнения
прогонов, агрегация по seed'ам; реализованы, но пока не прогнаны расширения --- ординальная
голова CORN, веса классов, групповой сплит, отбор эпохи по валидации, температурная
калибровка. Всего посчитано 12 прогонов дообучения ($\approx 2.5$ часа на Mac M4)
и 10 пробингов (41 минута).

Главные выводы: (1) stance воспроизводится в пределах разброса, premise --- на 4--6 пунктов
ниже авторского; (2) более половины преимущества rubert-base на stance объясняется явной
оценкой в тексте, а ruRoberta на ней не держится вовсе; (3) из-за этого исходная постановка
занижает разницу между энкодерами примерно вчетверо; (4) разброс по seed'ам и даже
по перезапуску с тем же seed'ом сопоставим с измеряемыми эффектами, поэтому одиночное
число в этой задаче не интерпретируемо.

Осталось (код готов, нужны вычислительные ресурсы уровня 4090): пять seed'ов для
ruRoberta-large в обоих режимах ($\approx 6$ часов на Mac); нижние ступени лестницы
дообучением, а не только пробингом; групповой сплит по фильму как отдельный режим оценки;
сравнение ординальной головы и взвешенного лосса с бейзлайном --- последнее имеет смысл
только на нескольких seed'ах, поскольку разброс $\pm 2.43$ съест любой эффект меньше
пяти пунктов.

\bibliographystyle{apalike}
\bibliography{lit}
\end{document}
